\chapter{\texorpdfstring{Vulnerability 5.39: Unprototyped\\Function Pointer (Stack OOB Write)}{Vulnerability 5.39: Unprototyped Function Pointer (Stack OOB Write)}}

\section{Executive Summary}

Vulnerability 5.39 is a \textbf{confirmed verifier bypass} combining an unprototyped function pointer (ISO-IEC Rule 5.39 violation) with a \textbf{tainted array index derived from the TCP sequence number}. The verifier loses range tracking for the tainted argument, but extended testing revealed that a \textbf{second, independent verifier check} on stack pointer arithmetic (\texttt{fp + register}) prevents the actual OOB write at runtime. The original code loads only because \textbf{LLVM eliminates the store as dead code} before the verifier sees it.

\begin{tcolorbox}[colback=blue!5!white,colframe=blue!75!black,title=ISO Rule Violated,fonttitle=\bfseries]
ISO-IEC Rule 5.39: ``Do not create pointers to functions with an implicit return type''
\end{tcolorbox}

\begin{tcolorbox}[colback=red!5!white,colframe=red!75!black,title=Severity,fonttitle=\bfseries]
\textbf{MEDIUM:} Verifier bypass confirmed (unprototyped function pointer breaks taint/range tracking). However, the actual OOB stack write is \textbf{independently blocked} by a second verifier check on \texttt{fp + register} arithmetic. Original code loads only because LLVM eliminates the store as dead code. \textbf{No runtime memory corruption in this scenario}.
\end{tcolorbox}

\subsection{Why This Vulnerability Matters}

\begin{enumerate}[label=\arabic*.]
  \item \textbf{Unprototyped Function Pointer}: Verifier loses tracking of what the function does and what arguments it accepts
  \item \textbf{Tainted Array Index}: TCP sequence number is attacker-controlled and used directly as an array index
  \item \textbf{Verifier Bypass}: Verifier cannot prove the index stays within bounds [0-15] because range info is lost at the unprototyped call boundary
  \item \textbf{Dead Store Elimination}: LLVM optimizes away the actual memory write, preventing the verifier from seeing the unsafe \texttt{fp + register} operation
  \item \textbf{Second Verifier Check}: When the write is forced with \texttt{volatile}, the verifier independently catches \texttt{fp + unbounded register}
  \item \textbf{Sandbox Constraint}: Even if the write occurred, eBPF execution is isolated from kernel memory
\end{enumerate}

\section{The Vulnerable Code Pattern}

\subsection{Vulnerable Code (Simplified)}

\vspace{0.3cm}
\noindent
\fcolorbox{black!30}{gray!15}{%
  \begin{minipage}[c]{\dimexpr\textwidth-2\fboxsep-2\fboxrule}
    {\small\ttfamily
    /* Target function with range expectation */\\
    struct sample\_struct \{\\
    \phantom{struct sample\_struct \{} \_\_u8 array1[16];  /* Expected 0-15 only */\\
    \phantom{struct sample\_struct \{} \_\_u8 array2[16];\\
    \};\\
    \\
    static \_\_always\_inline void restricted\_sink(int i) \{\\
    \phantom{static \_\_always\_inline void restricted\_sink(int i) \{} struct sample\_struct s = \{\};\\
    \phantom{static \_\_always\_inline void restricted\_sink(int i) \{} bpf\_printk("index : \%d", i);\\
    \phantom{static \_\_always\_inline void restricted\_sink(int i) \{} s.array1[i] = 42;  /* OUT-OF-BOUNDS if i >= 16 */\\
    \}\\
    \\
    /* Unprototyped function pointer */\\
    void (*pf)() = restricted\_sink;  /* No prototype! */\\
    \\
    /* Tainted input from TCP packet */\\
    \_\_u32 seq\_value = bpf\_htons(hdr-\textgreater{}tcp-\textgreater{}seq) / 1000;\\
    (*pf)(seq\_value);  /* Verifier loses range info at call boundary */
    }
  \end{minipage}%
}
\vspace{0.3cm}

\subsection{Why This Is Vulnerable}

\begin{enumerate}[label=\arabic*.]
  \item \textbf{Unprototyped pointer}: `void (*pf)()`, No parameter type info, verifier can't track argument constraints
  \item \textbf{Tainted input}: TCP sequence number is attacker-controlled (values 0-$2^{32}-1$)
  \item \textbf{Division doesn't constrain}: `seq / 1000` is still unbounded (0-4,294,967)
  \item \textbf{Array size is 16}: Accessing array1[16+] writes beyond buffer into adjacent stack memory
  \item \textbf{Verifier failure}: Without knowing the expected range [0-15], verifier can't prevent OOB access
\end{enumerate}

\section{Multi-Stage Analysis: Evidence Chain}

\subsection{Stage 1: Source Code (Patch Location)}

The vulnerability is introduced via a Git patch:
\begin{tcolorbox}[colback=gray!10, colframe=gray!10, arc=2pt, boxrule=0pt]
\small\ttfamily
\texttt{XDPs/xdp\_synproxy/patches/5\_39\_taintnoproto/0001-feat-5.39-taintnoproto*.patch}
\end{tcolorbox}

\subsection{Stage 2: Bytecode (llvm-objdump)}

Using \texttt{llvm-objdump -S xdp\_synproxy\_kern.o}:

\vspace{0.3cm}
\noindent
\fcolorbox{black!30}{gray!15}{%
  \begin{minipage}[c]{\dimexpr\textwidth-2\fboxsep-2\fboxrule}
    {\small\ttfamily
    ;~ (*pf)(bpf\_htons(hdr-\textgreater{}tcp-\textgreater{}seq)/1000);\\
    139: (61) r3 = *(u32 *)(r8 +4)       \# Load TCP seq\\
    140: (dc) r3 = be16 r3                \# Byte-order swap\\
    141: (37) r3 /= 1000                  \# Division (unbounded)\\
    ;~ bpf\_printk("index : \%d", i);\\
    142: (18) r1 = 0x18 ll\\
    145: (85) call 0x6                    \# Call with tainted index in R3\\
    ;~ s.array1[i] = 42;\\
    149: (73) *(u8 *)(r10 -80) = r3      \# STORE at offset -80 (UNSAFE!)
    }
  \end{minipage}%
}
\vspace{0.3cm}

\noindent
\textbf{Interpretation}:
- R3 is loaded from packet (line 139-141)
- No bounds check on R3
- Store instruction uses R3 as an offset (line 149)
- If R3 $\geq$ 16, this writes beyond the array buffer

\noindent
\textbf{Finding}: Bytecode shows the unprototyped call with tainted index. Not optimized away.

\subsection{Stage 3: Verifier (bpftool prog load)}

Command:
\begin{tcolorbox}[colback=gray!10, colframe=gray!10, arc=2pt, boxrule=0pt]
\small\ttfamily
\texttt{sudo bpftool prog load xdp\_synproxy\_kern.bpf.o /sys/fs/bpf/test\_synproxy type xdp -d}
\end{tcolorbox}

\noindent
\textbf{Result}: The verifier \textbf{accepts the program without error}.

\noindent
\textbf{Finding}: No rejection for the tainted array index. The unprototyped function call breaks taint tracking.

\subsection{Stage 4: Xlated Bytecode (bpftool prog dump xlated)}

\vspace{0.3cm}
\noindent
\fcolorbox{black!30}{gray!15}{%
  \begin{minipage}[c]{\dimexpr\textwidth-2\fboxsep-2\fboxrule}
    {\small\ttfamily
    139: (61) r3 = *(u32 *)(r8 +4)\\
    140: (dc) r3 = be16 r3\\
    141: (37) r3 /= 1000\\
    149: (73) *(u8 *)(r10 -80) = r3     \# Unsafe store still present\\
    }
  \end{minipage}%
}
\vspace{0.3cm}

\noindent
\textbf{Finding}: The xlated bytecode is nearly identical to compiled bytecode. \textbf{No bounds check inserted}. The vulnerable store survives verification unchanged.

\subsection{Stage 5: Runtime Execution (trace\_pipe)}

Trigger the vulnerability:

\begin{tcolorbox}[colback=gray!10, colframe=gray!10, arc=2pt, boxrule=0pt]
\small\ttfamily
python3 src/exploits/5\_39\_taintnoproto/poc\_539\_exploit.py 
\end{tcolorbox}

\noindent
Kernel trace output:

\vspace{0.3cm}
\noindent
\fcolorbox{black!30}{gray!15}{%
  \begin{minipage}[c]{\dimexpr\textwidth-2\fboxsep-2\fboxrule}
    {\small\ttfamily
    <idle>-0 [000] ..s2. 348.737756: bpf\_trace\_printk: [5.39]: index: 16\\
    <idle>-0 [000] ..s2. 348.738659: bpf\_trace\_printk: [5.39]: Content: 42\\
    <idle>-0 [000] ..s2. 348.739616: bpf\_trace\_printk: [5.39]: index: 17\\
    <idle>-0 [000] ..s2. 348.740573: bpf\_trace\_printk: [5.39]: Content: 42\\
    <idle>-0 [000] ..s2. 348.741530: bpf\_trace\_printk: [5.39]: index: 25\\
    <idle>-0 [000] ..s2. 348.742487: bpf\_trace\_printk: [5.39]: Content: 42
    }
  \end{minipage}%
}
\vspace{0.3cm}

\noindent
\textbf{Finding}: The trace shows indices 16, 17, 25 (all $\geq$ 16) and the value 42 printed at each index. However, this output requires careful interpretation: the \texttt{bpf\_printk("Content \%d", s.array1[i])} prints 42 because it is a hardcoded constant, not because it reads from a successfully written OOB location. LLVM eliminated the actual \texttt{s.array1[i] = 42} store as a dead store optimization before the bytecode reached the verifier.

\section{Extended Testing: Volatile Modification}

To verify whether the OOB write actually occurs at runtime (rather than just the index computation), a modified patch was created using \texttt{volatile} to force the compiler to emit the store instruction.

\subsection{The Test}

\noindent
The \texttt{sample\_struct} was declared as \texttt{volatile} and reads of adjacent array elements were added after the write, preventing LLVM from eliminating the store as dead code.

\subsection{Result: REJECTED by the Verifier}

\begin{tcolorbox}[colback=gray!10, colframe=gray!10, arc=2pt, boxrule=0pt]
\small\ttfamily
math between fp pointer and register with unbounded min value is not allowed
\end{tcolorbox}

\subsection{Interpretation}

\noindent
The verifier has \textbf{two independent safety checks}:

\begin{table}[h]
\centering
\begin{tabular}{|c|c|c|}
\hline
\textbf{Check} & \textbf{Purpose} & \textbf{5.39 Result} \\
\hline
Taint tracking (prototype) & Track range through calls & BYPASSED \\
Stack pointer arithmetic & Ensure fp+register bounded & NOT BYPASSED \\
\hline
\end{tabular}
\caption{Dual-layer verifier protection for vulnerability 5.39}
\end{table}

\noindent
\textbf{Why the original code loads}: LLVM sees that \texttt{s.array1[i] = 42} is never read after the write, eliminates the store as dead code, and without the store, the verifier never sees \texttt{fp + unbounded\_register}.

\noindent
\textbf{Why the volatile version is rejected}: \texttt{volatile} prevents LLVM from eliminating the store, so the bytecode contains \texttt{fp + r\_tainted}, and the verifier catches the unbounded register.

\subsection{Conclusion from Extended Testing}

\noindent
The verifier bypass via unprototyped function pointer is \textbf{real and confirmed}: the verifier loses range tracking for the tainted parameter. However, the practical impact is limited: without \texttt{volatile}, the compiler eliminates the OOB write; with \texttt{volatile}, the verifier's second check catches it. This creates a paradox where the vulnerability is only loadable when the dangerous operation does not actually execute.

\section{Exploitation Analysis}
\noindent
The verifier bypass has been \textbf{confirmed}: the unprototyped function pointer causes the verifier to lose range tracking for tainted arguments. The tainted index (TCP sequence number) reaches the function with out-of-bounds values (16, 17, 25), as confirmed by \texttt{bpf\_trace\_printk} output.

\vspace{0.2cm}
\noindent
However, extended testing with \texttt{volatile} revealed that the \textbf{actual memory write is prevented} by a dual-layer protection mechanism: (1) without \texttt{volatile}, LLVM eliminates the store as dead code before the verifier sees it; (2) with \texttt{volatile}, the verifier's independent check on \texttt{fp + register} arithmetic catches the unbounded register and rejects the program.

\vspace{0.2cm}
\noindent
The vulnerability demonstrates a \textbf{genuine verifier deficiency}: the unprototyped function pointer breaks taint/range tracking, which is a real gap in the verifier's safety analysis. However, the \textbf{defense-in-depth} of the verifier (the separate \texttt{fp + register} bounds check) independently prevents the actual OOB write from occurring at runtime.

\vspace{0.2cm}
\noindent
This finding is significant because it shows that the verifier's safety depends on \textbf{multiple independent checks working together}: if only the taint tracking check existed, the vulnerability would be exploitable whenever the compiler preserves the store. The second check provides a safety net, but the underlying verifier gap (loss of range info through unprototyped calls) remains a real concern for future code patterns.

\section{Artifact Evidence}

All analysis logs available in: \texttt{src/exploits/5\_39\_taintnoproto/logs/}

\begin{itemize}
  \item \texttt{5\_39\_bytecode\_analysis.txt}, Full llvm-objdump output
  \item \texttt{5\_39\_verifier\_log.txt}, Verifier output (bpftool prog load -d)
  \item \texttt{5\_39\_xlated\_dump.txt}, Xlated bytecode (verifier-processed)
\end{itemize}

\noindent
PoC script: \texttt{src/exploits/5\_39\_taintnoproto/poc\_539\_exploit.py}

